\section{Discussion and Limitations}
\label{sec:discussion}

\paragraph{When the interval can be shortened.}
Holding the labeled set, the unlabeled set, and the predictor fixed, a shorter valid interval requires three things at once: the calibration function $m$ must be nonlinear in the information used, the nonlinear part must be material compared with the linear residual variance ($\kappa$ large), and the labeled sample must be large enough that a flexible calibrator's error is smaller than that nonlinear part. The first two are properties of the problem; the third depends on $n$. Our simulations reach the bound up to the learner's error. The public pools satisfy the first condition only weakly: their nonlinear shares are at most $5\%$, which bounds the available gain at roughly $2\%$ of the interval width at $N=4n$.

\paragraph{Practical recipe.}
We recommend the following. (i) Fit a linear calibrator and compute $\kappa$ on the labeled data with cross-fitting; if it is small, report PPI++ and stop. (ii) If $\kappa$ is large, fit a smooth calibrator with cross-fitting and check whether its out-of-fold error is below $\Delta_L$; if it is not, the linear interval is the better choice. (iii) Report the chosen estimator and its alternatives, since the choice is data-dependent. The formulas in Theorem~\ref{thm:bound} and Corollary~\ref{cor:learned} give the bound and the break-even condition, and both are computable from the labeled sample.

\paragraph{Why more flexibility can hurt.}
Cross-fitting protects validity: the estimator stays unbiased for any learner. It does not protect width. A learner adds $(1-a)\mathrm{Var}(\hat m-m)/n$ to the variance, which is the reason a spline calibrator widens intervals by $16\%$ when $m$ is linear and $n=50$. Boosting, which has more variance at small $n$, widens them by $19\%$. In practice, a cross-validated choice between linear and flexible calibrators would reduce this cost, but it adds a selection step whose effect on coverage we did not study.

\paragraph{Comparison with stratification.}
Stratified PPI \citep{fisch2024strat} uses a piecewise-constant calibrator on the predictor. In our experiments it did not gain over PPI++ in any scenario, and it was $1$--$8\%$ wider on the real pools. We note a point that matters for implementation: when strata weights are estimated from the unlabeled data, as is standard, the variance must include the between-stratum term $\mathrm{Var}_U(f+\bar r_{b(f)})$. Treating the weights as known, as the variance formula in that paper does, undercovers in our development runs (coverage 0.86 across development cells, and a width ratio of about 0.67 on the linear scenario). The corrected linearization is the one used in this paper.

\paragraph{Limitations.}
\begin{itemize}
  \item \emph{Asymptotic statements.} Propositions~\ref{prop:var} and Theorem~\ref{thm:bound} are first-order. The variance formula is a Wald approximation; at $n=50$ with heavy-tailed or discrete outcomes it undercovers (census income) or over-covers (forest). Exact or bounded-outcome intervals could be tighter at small $n$; we did not evaluate them.
  \item \emph{Mean target and i.i.d. sampling.} We study $\mathbb E[Y]$, which includes prevalence. Quantiles, regression coefficients, and covariate or label shift are not covered. The real-data target is the pool mean, and the pools are treated as populations with i.i.d.\ resampling.
  \item \emph{Real data coverage.} Only seven pools had full labels and predictions. \texttt{plankton} has labels on a different year than its unlabeled set, which is a distribution-shift problem outside our question, and \texttt{ballots} has no labels on its unlabeled pool, so coverage is undefined. We did not run folktables or other external datasets.
  \item \emph{Estimated bound.} The pool-level bound uses a cross-fitted gradient-boosting estimate of $m$. For some pools that estimate is less accurate than the linear fit, so the reported bound is conservative: $\Delta_L$ is inflated by estimation error. The bound columns are therefore an upper estimate of the available gain.
  \item \emph{Learner choice.} The spline calibrator was selected on synthetic development data. Other learners could do better on particular real pools; Table~\ref{tab:realsens} reports all three we tried.
  \item \emph{Sample size.} Our main grid goes to $n=1000$; on four pools we also ran $n=5{,}000$ and $20{,}000$ (\cref{sec:results}). The realized gain approaches the bound only slowly: on \texttt{forest} at $n=20{,}000$ it is half a percent against a bound of two percent. Larger labeled samples than we could check may be needed before the flexible estimator reaches the bound on pools with small $\kappa$.
\end{itemize}

\paragraph{Scope of the claim.}
The result concerns regular Wald-type intervals that use the same $(Z_L,Y_L,Z_U)$. It does not rule out shorter intervals from different assumptions, such as bounded outcomes or additional unlabeled covariates beyond those we used. Within the stated information set, the efficient bound is a lower bound on asymptotic variance for regular estimators.
